\begin{definition} \label{an1:def:connected-set}
  Sets $A,B\subseteq X$ are \textbf{separated} if $\overline A\cap B=A\cap\overline B=\varnothing$.
  A set $E$ is \textbf{connected} if it cannot be expressed as the union of two nonempty separated sets; otherwise it is \textbf{disconnected}.
  The empty set is connected under this convention. The union equality automatically makes both sets subsets of $E$.
  \begin{lean}{Two sets are separated when neither meets the closure of the other.}
    def separated {X : Type*} [MetricSpace X] (A B : Set X) : Prop :=
      Disjoint (closure A) B ∧ Disjoint A (closure B)
  \end{lean}
  \begin{lean}{A connected set admits no separation into two nonempty parts.}
    def connected {X : Type*} [MetricSpace X] (E : Set X) : Prop :=
      ∀ A B : Set X, A.Nonempty → B.Nonempty → E = A ∪ B → ¬ separated A B
  \end{lean}
\end{definition}

% Lean source: Textbooks.MathematicalAnalysis1.Chapter01.ConnectedSets.closure_upper_bound
\begin{lemma} \label{an1:lem:closure-upper-bound}
  \begin{lean}{If every point of $A\subseteq\R$ is at most $c$, the same is true for every point of $\overline A$.}
    theorem closure_upper_bound (A : Set ℝ) (c : ℝ) (h : ∀ x ∈ A, x ≤ c) :
        ∀ x ∈ closure A, x ≤ c
  \end{lean}
\end{lemma}

\begin{proof}
  \begin{lean}{If $x\in\overline A$ were greater than $c$, choose a point $y\in A$ within $x-c$ of $x$. Then $|x-y|<x-c$ would imply $y>c$, contradicting the upper bound.}
    intro x hx
    by_contra hn
    have hc : c < x := lt_of_not_ge hn
    obtain ⟨y, hy, hd⟩ := Metric.mem_closure_iff.mp hx (x - c) (sub_pos.mpr hc)
    rw [Real.dist_eq, abs_lt] at hd
    have hyc := h y hy
    linarith
  \end{lean}

\end{proof}

% Lean source: Textbooks.MathematicalAnalysis1.Chapter01.ConnectedSets.closure_lower_bound
\begin{lemma} \label{an1:lem:closure-lower-bound}
  \begin{lean}{If every point of $A\subseteq\R$ is at least $c$, the same is true for every point of $\overline A$.}
    theorem closure_lower_bound (A : Set ℝ) (c : ℝ) (h : ∀ x ∈ A, c ≤ x) :
        ∀ x ∈ closure A, c ≤ x
  \end{lean}
\end{lemma}

\begin{proof}
  \begin{lean}{If $x\in\overline A$ were less than $c$, a point $y\in A$ within $c-x$ of $x$ would satisfy $y<c$. This contradicts the lower bound.}
    intro x hx
    by_contra hn
    have hc : x < c := lt_of_not_ge hn
    obtain ⟨y, hy, hd⟩ := Metric.mem_closure_iff.mp hx (c - x) (sub_pos.mpr hc)
    rw [Real.dist_eq, abs_lt] at hd
    have hcy := h y hy
    linarith
  \end{lean}

\end{proof}

% Lean source: Textbooks.MathematicalAnalysis1.Chapter01.ConnectedSets.no_separation_of_between
\begin{lemma} \label{an1:lem:no-separation-of-between}
  \begin{lean}{Suppose $E\subseteq\R$ contains every point strictly between any two of its elements.
It is impossible to have $E=A\cup B$ with $A,B$ separated and points $a\in A$, $b\in B$, $a<b$.}
    theorem no_separation_of_between (E A B : Set ℝ)
        (hbetween : ∀ a ∈ E, ∀ b ∈ E, ∀ x : ℝ, a < x → x < b → x ∈ E)
        (heq : E = A ∪ B) (hsep : separated A B)
        (a b : ℝ) (ha : a ∈ A) (hb : b ∈ B) (hab : a < b) : False
  \end{lean}
\end{lemma}

\begin{proof}
  \begin{lean}{Let $S=A\cap[a,b]$. It contains $a$ and is bounded above by $b$, so $c=\sup S$ exists and satisfies $a\leq c\leq b$. The supremum is in $\overline S\subseteq\overline A$.}
    let S := A ∩ Icc a b
    have hane : a ∈ S := ⟨ha, le_rfl, hab.le⟩
    have hne : S.Nonempty := ⟨a, hane⟩
    have hbd : BddAbove S := ⟨b, fun x hx => hx.2.2⟩
    let c := sSup S
    have hac : a ≤ c := le_csSup hbd hane
    have hcb : c ≤ b := csSup_le hne (fun x hx => hx.2.2)
    have hcS : c ∈ closure S := (supremum_in_closure S hne hbd).1
    have hcA : c ∈ closure A := (closure_properties S).2.2 (closure A)
      (closure_closed A) (fun x hx => subset_closure hx.1) hcS
  \end{lean}

  \begin{lean}{Separation implies $c\notin B$, and thus $c<b$. If $c=a$, it is already in $E$; otherwise it is strictly between $a$ and $b$, so the between-points hypothesis puts it in $E$. Since $c\notin B$, we have $c\in A$.}
    have hcnotB : c ∉ B := fun h => Set.disjoint_left.mp hsep.1 hcA h
    have hcltb : c < b := lt_of_le_of_ne hcb (fun h => hcnotB (h.symm ▸ hb))
    have hcE : c ∈ E := by
      by_cases he : a = c
      · rw [← he, heq]
        exact Or.inl ha
      · exact hbetween a (heq.symm ▸ Or.inl ha) b (heq.symm ▸ Or.inr hb) c
          (lt_of_le_of_ne hac he) hcltb
    have hcain : c ∈ A := (show c ∈ A ∪ B from heq ▸ hcE).resolve_right hcnotB
  \end{lean}

  \begin{lean}{Now separation gives $c\notin\overline B$. Consequently some $r>0$ satisfies $d(c,y)\geq r$ for all $y\in B$.}
    have hcnotclB : c ∉ closure B := fun h => Set.disjoint_left.mp hsep.2 hcain h
    have hex : ∃ r > 0, ∀ y ∈ B, r ≤ dist c y := by
      by_contra h
      push Not at h
      exact hcnotclB (Metric.mem_closure_iff.mpr h)
    obtain ⟨r, hr, hfar⟩ := hex
  \end{lean}

  \begin{lean}{Choose $\delta=\min(r,b-c)/2$. It is positive, less than both $r$ and $b-c$, and $x=c+\delta$ lies strictly between $a$ and $b$. Hence $x\in E$.}
    let δ := min r (b - c) / 2
    have hδ : 0 < δ := half_pos (lt_min hr (sub_pos.mpr hcltb))
    have hδr : δ < r := (half_lt_self (lt_min hr (sub_pos.mpr hcltb))).trans_le (min_le_left _ _)
    have hδb : δ < b - c := (half_lt_self (lt_min hr (sub_pos.mpr hcltb))).trans_le (min_le_right _ _)
    let x := c + δ
    have hax : a < x := by dsimp [x]; linarith
    have hxb : x < b := by dsimp [x]; linarith
    have hxE := hbetween a (heq.symm ▸ Or.inl ha) b (heq.symm ▸ Or.inr hb) x hax hxb
  \end{lean}

  \begin{lean}{The distance $d(c,x)=\delta<r$ excludes $x$ from $B$. Therefore $x\in A\cap[a,b]=S$, yet $x>c=\sup S$, a contradiction. Choosing half the available radius ensures that the new point is inside the excluded ball, not merely on its boundary.}
    have hxnotB : x ∉ B := by
      intro hx
      have hdist := hfar x hx
      have hd : dist c x = δ := by
        rw [Real.dist_eq, show c - x = -δ by dsimp [x]; ring, abs_neg, abs_of_pos hδ]
      rw [hd] at hdist
      linarith
    have hxA := (show x ∈ A ∪ B from heq ▸ hxE).resolve_right hxnotB
    have hxc : x ≤ c := le_csSup hbd (show x ∈ S from ⟨hxA, hax.le, hxb.le⟩)
    dsimp [x] at hxc
    linarith
  \end{lean}

\end{proof}

% Lean source: Textbooks.MathematicalAnalysis1.Chapter01.ConnectedSets.connected_iff_between
\begin{theorem} \label{an1:thm:connected-sets-in-R}
  \begin{lean}{A set $E\subseteq\R$ is connected if and only if, whenever $a,b\in E$ and $a<x<b$, one has $x\in E$.}
    theorem connected_iff_between (E : Set ℝ) :
        connected E ↔ ∀ a ∈ E, ∀ b ∈ E, ∀ x : ℝ, a < x → x < b → x ∈ E
  \end{lean}
\end{theorem}

\begin{proof}
  \begin{lean}{Suppose $E$ is connected and a point $x$ between $a,b\in E$ were missing. Split $E$ into $A=E\cap(-\infty,x)$ and $B=E\cap(x,\infty)$. They cover $E$, and $a\in A$, $b\in B$ make both nonempty.}
    classical
    constructor
    · intro hc a ha b hb x hax hxb
      by_contra hx
      let A := E ∩ Iio x
      let B := E ∩ Ioi x
      have heq : E = A ∪ B := by
        ext y
        constructor
        · intro hy
          rcases lt_trichotomy y x with h | h | h
          · exact Or.inl ⟨hy, h⟩
          · exact False.elim (hx (h ▸ hy))
          · exact Or.inr ⟨hy, h⟩
        · rintro (h | h) <;> exact h.1
  \end{lean}

  \begin{lean}{The closure of $A$ stays at or below $x$, so it misses $B$. Similarly the closure of $B$ stays at or above $x$, so it misses $A$. This is a forbidden separation.}
    apply hc A B ⟨a, ha, hax⟩ ⟨b, hb, hxb⟩ heq
    constructor
    · apply Set.disjoint_left.mpr
      intro y hy hB
      have hle := closure_upper_bound A x (fun z hz => hz.2.le) y hy
      exact (not_lt_of_ge hle) hB.2
    · apply Set.disjoint_left.mpr
      intro y hA hy
      have hle := closure_lower_bound B x (fun z hz => hz.2.le) y hy
      exact (not_lt_of_ge hle) hA.2
  \end{lean}

  \begin{lean}{Conversely, suppose every intermediate point is in $E$ and assume a separation $E=A\cup B$. Choose $a\in A$ and $b\in B$; separation makes them distinct. If $a<b$, the preceding supremum argument contradicts the separation. If $b<a$, interchange the two sets and apply the same argument.}
    · intro h A B hA hB heq hsep
      obtain ⟨a, ha⟩ := hA
      obtain ⟨b, hb⟩ := hB
      have hne : a ≠ b := by
        intro he
        exact Set.disjoint_left.mp hsep.1 (subset_closure ha) (he.symm ▸ hb)
      rcases lt_or_gt_of_ne hne with hab | hba
      · exact no_separation_of_between E A B h heq hsep a b ha hb hab
      · have hsep' : separated B A := ⟨hsep.2.symm, hsep.1.symm⟩
        exact no_separation_of_between E B A h (heq.trans (union_comm _ _)) hsep' b a hb ha hba
  \end{lean}

\end{proof}
